\section{Regularity and existence of the elliptic deformation}
\label{n3:existence:section}

We now pass from the bounded variables in
Proposition~\ref{n3:apriori:bounds} to smooth solutions. The second equation
contains Hessian squares, so classical estimates cannot yet be applied
directly. We first control the gradient of the first unknown with a
measurable axial coefficient. Its Hessian Morrey bound then gives
continuity of the second unknown and permits classical regularity.
Finally we define the scalar Dirichlet solution operator on all bounded
input sets and apply degree to
Proposition~\ref{n3:deformation:homotopy}.

Throughout this section, ordinary regularity constants may depend on
all fixed deformation parameters, including $n$ and $\tau$.  No such
constant will be used as a polynomial-in-$n$ estimate.  Bounds uniform
in the outer truncation radius will be stated explicitly.

\subsection{An estimate for a measurable axial coefficient}

We use $D$ for coordinate derivatives in the analytic lemmas.  Balls
are Euclidean coordinate balls, and a boundary ball is their
intersection with one side of a smooth boundary.  A family of boundary
patches has bounded geometry here if normal-coordinate charts, their
inverses, the metric and its inverse, and the derivatives through the
orders used below have fixed bounds.  We only apply the lemmas to
smooth metrics and smooth boundaries.  Their doubled metrics will be
Lipschitz and smooth on each side of a single plane.

\begin{lemma}[Axial gradient estimate]\label{n3:existence:gradient}
Let $U'$ be compactly contained in an interior or boundary coordinate
patch $U$ in dimension three.  In a boundary patch assume that $z$ is
constant on the boundary face.  Suppose $z$ is smooth up to that face,
$|\nabla z|\le L$, and, almost everywhere,
\begin{equation}\label{n3:existence:axial-equation}
 A_\chi:\Hess z=G,\qquad
 A_\chi=I-(1-\chi)e\otimes e,\qquad
 e=\frac{\nabla z}{|\nabla z|},\qquad
 0<\chi_0\le\chi\le1,
\end{equation}
where $\|G\|_\infty\le Q$.  At zero gradients any measurable unit
axis for which the equation holds is allowed.  Then there are
$\alpha\in(0,1)$ and $C<\infty$, depending only on $\chi_0$, the
patch geometry and the separation of $U'$ from the other patch
boundaries, such that
\begin{equation}\label{n3:existence:gradient-estimate}
 \|\nabla z\|_{C^{0,\alpha}(U')}
 \le C(L+Q).
\end{equation}
For all sufficiently small coordinate balls centered in $U'$,
with intersection with the domain understood at a boundary, one also
has
\begin{equation}\label{n3:existence:hessian-morrey}
 \int_{B_r\cap U}|\Hess z|^2\dd V_g
 \le C(L+Q)^2r^{1+2\alpha}.
\end{equation}
In particular, the constants are independent of continuity moduli of
$\chi$ and $G$.
\end{lemma}

\begin{proof}
The proof combines a Cordes estimate, a divergence identity for the
gradient, and improvement of affine approximations. We first establish
these estimates, then handle boundary reflection and the iteration.

\paragraph{A Cordes estimate above exponent two.}
For a Euclidean unit vector $e$,
\begin{equation}\label{n3:existence:cordes-defect}
 \|I-A_\chi\|_F=1-\chi\le1-\chi_0.
\end{equation}
For a compactly supported smooth function $v$, Plancherel's theorem
gives $\|D^2v\|_2=\|\Delta v\|_2$.  The Calder\'on--Zygmund
operator taking $\Delta v$ to the full array $D^2v$ is bounded on
$L^p$, and interpolation makes its norm arbitrarily close to one as
$p\to2$.  Choose once and for all
\begin{equation}\label{n3:existence:p-zero}
 2<p_0<3,\qquad
 C_{p_0}(1-\chi_0/2)<1.
\end{equation}
Rescale a metric-normalized patch until the coordinate principal
matrix differs from its Euclidean axial matrix by less than
$\chi_0/2$ in Frobenius norm.  Absorbing this difference in the
Laplacian estimate, applying a cutoff, and using
$\|Dv\|_p\le\eta\|D^2v\|_p+C_\eta\|v\|_p$ gives
\begin{equation}\label{n3:existence:local-cordes}
 \|D^2v\|_{L^{p_0}(B_{1/2})}
 +\|Dv\|_{L^{p_0}(B_{1/2})}
 \le C\bigl(\|v\|_{L^{p_0}(B_1)}
       +\|A_\chi:\Hess v\|_{L^{p_0}(B_1)}\bigr).
\end{equation}
The Christoffel terms are bounded first-order terms in this estimate.
The same argument applies when $v=z-L_0$ and $L_0$ is coordinate
affine: its extra right side is bounded by
$C\|Dg\|_\infty|DL_0|$.
The estimate uses bounded measurable principal coefficients only;
the smoothness of the solutions allows all of its applications
without an approximation issue.  The linear estimates just used are
the ordinary Laplacian estimates; see
\cite[Chapters 7 and 9]{GilbargTrudinger2001}.

\paragraph{The identity controlling a near-maximal gradient.}
Write $H=\Hess z$ and $P=\nabla z$.  Equation~\eqref{n3:existence:axial-equation}
gives the pointwise vector identity
\begin{equation}\label{n3:existence:gradient-flux}
 A_\chi\nabla\frac{|P|^2}{2}-GP
       =(H-\Delta z\,g)P.
\end{equation}
Indeed $HP=|P|H(e,\cdot)$ and
$G=\Delta z-(1-\chi)H(e,e)$, so the two axial terms cancel.
The identity is also valid at $P=0$.  Taking the divergence of its
right side, rather than differentiating $\chi$, yields
\begin{equation}\label{n3:existence:bochner-flux}
 \Div\left(A_\chi\nabla\frac{|P|^2}{2}-GP\right)
 =|H|^2-(\Delta z)^2+\Ric(P,P).
\end{equation}
Choose $a_0>0$ so small that
$(1+a_0)(1-\chi_0)^2<1$.  Young's inequality and the equation show
\begin{equation}\label{n3:existence:positive-hessian}
 |H|^2-(\Delta z)^2
 \ge c_0|H|^2-C_0G^2,
\end{equation}
where $c_0>0$ depends only on $\chi_0$.

Consider solutions normalized by $|\nabla z|\le1$ on $B_1$.
After making the scale small, let the forcing and the rescaled
geometric errors have size at most $\delta$.  The nonnegative function
$s=1-|\nabla z|^2$ then satisfies
\begin{equation}\label{n3:existence:s-supersolution}
 \Div(A_\chi\nabla s+2G\nabla z)
 \le -2c_0|\Hess z|^2+C\delta
\end{equation}
in a weak sense.  In the Euclidean case the last term can be replaced
by $C\delta^2$.  We allow $C\delta$ to include the geometric errors.

For the compactness step, suppose
the normalized errors tend to zero in a sequence and
$\inf_{B_{1/2}}s\to0$.  The inhomogeneous weak Harnack inequality,
applied to Equation~\eqref{n3:existence:s-supersolution} after dropping
its favorable Hessian term, gives $s\to0$ in measure on $B_{1/2}$.
The boundedness $0\le s\le1$ also gives convergence in every finite
$L^p$ norm there.  To recover the Hessian information, take a cutoff
$\eta$ supported inside $B_{1/2}$ and use
$\eta^2(b-s)_+$ in the weak supersolution inequality.  Ellipticity
and Young's inequality give
\begin{equation}\label{n3:existence:low-truncation}
 \int_{\{s<b\}}\eta^2|Ds|^2
 \le Cb^2\int|D\eta|^2+o(1)
\end{equation}
for each fixed $b>0$ as the errors vanish.  The error fields in
Equation~\eqref{n3:existence:s-supersolution} are bounded and tend to
zero, so their products with $Ds$ are absorbed in deriving this
inequality.  Equation~\eqref{n3:existence:local-cordes} supplies a
uniform $L^2$ bound for $Ds$.  Its integral over $\{s\ge b\}$ tends
to zero in $L^1$ by Cauchy's inequality and convergence in measure.
On $\{s<b\}$ use Equation~\eqref{n3:existence:low-truncation}, then let
$b\downarrow0$.  Thus $Ds\to0$ locally in $L^1$.
Testing Equation~\eqref{n3:existence:bochner-flux} with a nonnegative
compact cutoff and using Equations~\eqref{n3:existence:positive-hessian}
and \eqref{n3:existence:gradient-flux} now gives
\begin{equation}\label{n3:existence:hessian-flat-limit}
 \int_{B_{1/4}}|\Hess z|^2\longrightarrow0.
\end{equation}
Uniform Lipschitz compactness implies that, after subtracting a
constant and taking a subsequence, $z$ tends uniformly on a smaller
ball to an affine function.  Its slope has length one because
$|\nabla z|\to1$ in measure and the Hessians tend to zero.

Consequently, for any prescribed sufficiently small $b_*>0$ there
are fixed $r_*\in(0,1/4)$, $k_*\in(r_*,1)$ and $\delta_*>0$ such
that every normalized solution with errors at most $\delta_*$ has
one of the following properties:
\begin{equation}\label{n3:existence:drop-or-flat}
 \begin{split}
 &\sup_{B_{r_*}}|\nabla z|\le k_*;\qquad\text{or}\\
 &\sup_{B_{r_*}}|z-L_0|\le b_*r_*,
       \qquad \tfrac12\le|DL_0|\le2.
 \end{split}
\end{equation}
If this conclusion failed, solutions whose gradient suprema approach
one and whose errors approach zero would contradict
Equation~\eqref{n3:existence:hessian-flat-limit}.  A gradient-drop
constant can always be weakened toward one to ensure $k_*>r_*$.

\paragraph{A fixed-axis equation.}
Once a nonzero affine slope appears, we compare with an equation whose
axis is fixed. Fix a Euclidean unit vector $e_0$. Suppose
\begin{equation}\label{n3:existence:fixed-axis}
 \Delta_{e_0^\perp}Y_0+\chi(x)\partial_{e_0e_0}Y_0=0
 \quad\text{in }B_{3/4},\qquad
 \|Y_0\|_{W^{2,2}(B_{3/4})}\le C_1.
\end{equation}
The derivative $v=\partial_{e_0}Y_0$ belongs to $W^{1,2}$ and,
in distributions, satisfies
\begin{equation}\label{n3:existence:axial-derivative-divergence}
 \Delta_{e_0^\perp}v+
 \partial_{e_0}(\chi\,\partial_{e_0}v)=0.
\end{equation}
This follows by differentiating the distribution $\chi\partial_{e_0}v$
as a whole.  It does not assume a derivative of $\chi$ exists as a
function.  De Giorgi's estimate for this uniformly elliptic
divergence equation, followed by Caccioppoli applied to
$v-v(x_0)$, gives an exponent $\alpha_1\in(0,1)$ and
\begin{equation}\label{n3:existence:axial-derivative-morrey}
 [v]_{C^{0,\alpha_1}(B_{1/2})}\le C,
 \qquad
 \int_{B_r(x_0)}|Dv|^2\le Cr^{1+2\alpha_1}
\end{equation}
for balls in a smaller fixed interior region.  These are the scalar
divergence estimates with ellipticity $\chi_0$; see
\cite[Chapter 8]{GilbargTrudinger2001}.

All components of $DY_0$, not only $v$, are controlled by this
estimate.  In fact,
\begin{equation}\label{n3:existence:poisson-source}
 \Delta Y_0=f_0,\qquad
 f_0=(1-\chi)\partial_{e_0}v,
 \qquad
 \int_{B_r(x_0)}|f_0|\le Cr^{2+\alpha_1}.
\end{equation}
The last inequality follows from Cauchy's inequality and
Equation~\eqref{n3:existence:axial-derivative-morrey}.  Take the Newton
potential of $f_0$ restricted to a slightly larger interior ball;
its difference from $Y_0$ is harmonic in that ball.  For two points
at distance $h$, the contribution to the difference of gradient
potentials from annuli of radius $r\le2h$ is bounded by
$C\sum r^{\alpha_1}\le Ch^{\alpha_1}$.  On the remaining annuli
the difference of gradient kernels is at most $Ch r^{-3}$, so their
contribution is
\begin{equation}\label{n3:existence:potential-annuli}
 C h\sum_{r\ge2h}r^{\alpha_1-1}
 \le Ch^{\alpha_1}.
\end{equation}
The harmonic remainder has its usual interior derivative estimates.
Therefore, decreasing the interior ball if needed,
\begin{equation}\label{n3:existence:model-gradient}
 \|Y_0\|_{C^{1,\alpha_1}(B_{1/4})}\le C.
\end{equation}

\paragraph{Improvement of a nonzero affine approximation.}
Suppose now
\begin{equation}\label{n3:existence:flat-input}
 \|z-L_0\|_{L^\infty(B_1)}\le b,
 \qquad \tfrac14\le|DL_0|\le4,
 \qquad |\nabla z|\le8.
\end{equation}
Let the forcing and the first-derivative metric errors be at most
$b\delta$.  For $Y=(z-L_0)/b$, Equation~\eqref{n3:existence:local-cordes}
on fixed nested balls gives
\begin{equation}\label{n3:existence:normalized-residual}
 \|Y\|_{W^{2,p_0}(B_{7/8})}\le C.
\end{equation}
In particular, this estimate does not require a bound for
$DY$ in the supremum norm.  Since $Dz=DL_0+bDY$, the variable axis
converges in measure to $e_0=DL_0/|DL_0|$ as $b\to0$.  More
quantitatively, the coefficient error $E$ caused by replacing this
axis by $e_0$ is bounded and satisfies
\begin{equation}\label{n3:existence:axis-error}
 \|E\|_{L^q(B_{7/8})}\le Cb^{\vartheta}+Cb\delta,
 \qquad q=\frac{2p_0}{p_0-2},
 \qquad \vartheta=\frac{p_0-2}{2}>0.
\end{equation}
Indeed the unit-direction map is Lipschitz where
$|bDY|<|DL_0|/2$, and elsewhere its difference is bounded; the
$L^{p_0}$ bound in Equation~\eqref{n3:existence:normalized-residual}
and interpolation give Equation~\eqref{n3:existence:axis-error}.
The same estimate includes the metric principal-part error.
H\"older's inequality, with $1/2=1/q+1/p_0$, shows that
\begin{equation}\label{n3:existence:small-fixed-axis-source}
 \bigl\|[I-(1-\chi)e_0\otimes e_0]:D^2Y\bigr\|_{L^2(B_{3/4})}
 \le C(b^{\vartheta}+\delta).
\end{equation}

Solve for a zero-Dirichlet correction $w$ on $B_{3/4}$ with the
left side of Equation~\eqref{n3:existence:small-fixed-axis-source} as
its source.  This does not require a nondivergence Green function.
The convex Dirichlet Hessian inequality
$\|D^2w\|_2\le\|\Delta w\|_2$ and
Equation~\eqref{n3:existence:cordes-defect} make the Laplace solution
operator a contraction on $W^{2,2}\cap W^{1,2}_0$.  Hence
\begin{equation}\label{n3:existence:fixed-axis-correction}
 \|w\|_{W^{2,2}(B_{3/4})}
 \le C\chi_0^{-1}(b^{\vartheta}+\delta),\qquad
 \|w\|_\infty\le C(b^{\vartheta}+\delta).
\end{equation}
The last implication uses $W^{2,2}\hookrightarrow C^{0,1/2}$ in
dimension three.  The convex Hessian inequality is the
Miranda--Talenti estimate; see \cite[Section~2, Theorem~3]{Talenti1965}.
Now $Y_0=Y-w$ satisfies Equation~\eqref{n3:existence:fixed-axis} with
a uniform $W^{2,2}$ bound.

Choose $0<\alpha_2<\alpha_1$ and then a fixed
$\lambda_0\in(0,1/4)$ such that
$C\lambda_0^{1+\alpha_1}\le\lambda_0^{1+\alpha_2}/4$ in
Equation~\eqref{n3:existence:model-gradient}.  By first taking
$\delta$ small and then $b\le b_*$ small,
Equation~\eqref{n3:existence:fixed-axis-correction} has supremum norm
at most $\lambda_0^{1+\alpha_2}/4$.  Taylor expansion of $Y_0$
at the center therefore gives an affine function $L_1$ with
\begin{equation}\label{n3:existence:improved-flatness}
 \sup_{B_{\lambda_0}}|z-L_1|
 \le b\lambda_0^{1+\alpha_2},\qquad
 |DL_1-DL_0|\le Cb.
\end{equation}
All constants are uniform in the slope range in
Equation~\eqref{n3:existence:flat-input}.

\paragraph{Reflection at a constant Dirichlet face.}
Subtract the boundary value and use normal coordinates with the face
given by $x_3=0$.  Reflect the metric across this plane and reflect
$z$ oddly.  Tangential derivatives of $z$ vanish on the face, so
the extension is $C^1$.  There is no distributional Hessian atom.
The doubled metric is Lipschitz; the coefficient and lower-order
estimates in Equation~\eqref{n3:existence:local-cordes} hold with its
essential first-derivative bound.

The flux in Equation~\eqref{n3:existence:gradient-flux} may have a
normal jump.  On either side of the face its normal trace is
\begin{equation}\label{n3:existence:reflected-flux}
 \bigl((\Hess z-\Delta z\,g)\nabla z\bigr)_\nu
   =-(\tr_T\Hess z)\,\partial_\nu z.
\end{equation}
Because the boundary value is constant, $\tr_T\Hess z$ is its
mean-curvature coefficient times $\partial_\nu z$.  The jump is
therefore bounded by $C|\mathrm{II}|\,|\nabla z|^2$.  At scale
$r$ it is a bounded plane density of size $O(r)$.  Write the plane
density and flux after multiplication by the reflected coordinate
volume density, so that the following divergence is Euclidean.
Such a signed density $a(x')\,\delta_{\{x_3=0\}}$ is the divergence of the
bounded field $a(x')\mathbf 1_{\{x_3>0\}}\partial_3$.
Thus its negative part is another small divergence error in
Equation~\eqref{n3:existence:s-supersolution}.  Smooth-side curvature
errors are $O(r^2)$.  Weak Harnack and the low-truncation argument
remain valid, and the flatness comparison uses the same Lipschitz
metric estimates.  This proves all the preceding alternatives for
balls meeting the reflection plane as well.

\paragraph{Iteration and the Morrey consequence.}
Keep $\alpha_2,\lambda_0$ from the affine improvement fixed.  Choose
$b_*$ sufficiently small that
$Cb_*/(1-\lambda_0^{\alpha_2})<1/4$, with $C$ the slope-increment
constant in Equation~\eqref{n3:existence:improved-flatness}.
Then choose $r_*,k_*,\delta_*$ in
Equation~\eqref{n3:existence:drop-or-flat} for this fixed $b_*$.
Normalize initially by a fixed multiple of $L+Q$ and choose a
uniformly small initial radius so that the metric and forcing errors
satisfy all the preceding smallness conditions.  If the first case
of Equation~\eqref{n3:existence:drop-or-flat} holds, rescale the ball
by $r_*$ and the gradient bound by $k_*$.  The normalized forcing
decreases by $r_*/k_*<1$.  Repeating this step either continues
forever or reaches the second case.  In the former case the gradient
decays with exponent $\log k_*/\log r_*$.
In the latter case apply Equation~\eqref{n3:existence:improved-flatness}
successively, retaining exponent $\alpha_2$.  At each affine step,
rescale $z$ by dividing by the spatial scale, subtracting only an
additive constant; this preserves its gradient bound.
Its slope changes have total size at most
$Cb_*/(1-\lambda_0^{\alpha_2})$, so the slopes remain in $[1/4,4]$.
At the $j$th step the error size is $b_*\lambda_0^{j\alpha_2}$,
while the normalized forcing and metric first-derivative errors have
size $O(\lambda_0^j)$.  Since $\alpha_2<1$, their ratio to the error
size only decreases.  This justifies every subsequent application
of the improvement lemma.  To combine the two branches, now choose
\begin{equation}\label{n3:existence:iteration-exponent}
 0<\alpha\le\min\left\{\alpha_2,
              \frac{\log k_*}{\log r_*}\right\},\qquad \alpha<1.
\end{equation}

At every center we have consequently obtained affine functions
$L_{x,r}$ satisfying
\begin{equation}\label{n3:existence:campanato-affine}
 \sup_{B_r(x)}|z-L_{x,r}|
 \le C(L+Q)r^{1+\alpha}.
\end{equation}
The slope differences at consecutive radii are bounded by
$C(L+Q)r^\alpha$; comparing the affine approximations on
overlapping balls gives Equation~\eqref{n3:existence:gradient-estimate}.
This also identifies their limiting slope with $Dz$, since $z$ is
smooth.  The reflected version gives the estimate up to a boundary.

Finally apply the $L^2$ version of
Equation~\eqref{n3:existence:local-cordes} to $z-L_{x,2r}$ after
rescaling.  The affine error contributes
$C(L+Q)r^{1/2+\alpha}$ to the Hessian $L^2$ norm; the bounded
forcing and affine Christoffel term contribute $C(L+Q)r^{3/2}$.
Squaring and using $\alpha<1$ proves
Equation~\eqref{n3:existence:hessian-morrey}.
\end{proof}

\subsection{Bounded solutions with quadratic gradient growth}

The Hessian Morrey bound controls the forcing in the second equation
through the following oscillation estimate. A positive measure in the
source bound also accommodates reflected boundary terms.

\begin{lemma}[Natural-growth oscillation estimate]
\label{n3:existence:natural-growth}
On a three-dimensional ball suppose $Z$ is bounded, $|Z|\le M$,
and is a smooth solution, or a reflected weak solution, of
\begin{equation}\label{n3:existence:natural-equation}
 \Div(aDZ+b)=e.
\end{equation}
Assume $a$ is measurable and symmetric,
$\lambda I\le a\le\Lambda I$, $|b|\le B$, and
\begin{equation}\label{n3:existence:natural-source}
 |e|\le C_0(1+|DZ|^2)\dd x+\mu,\qquad
 \mu(B_r(x))\le C_\mu r^{1+\beta}
\end{equation}
for some $\beta>0$ and every sufficiently small ball in the patch.
Then $Z$ has a local H\"older estimate whose constants depend only
on the displayed bounds and the patch separation.  The conclusion
also holds up to a smooth boundary with bounded conormal flux or
constant Dirichlet data, provided the corresponding reflected
measure satisfies Equation~\eqref{n3:existence:natural-source}.
\end{lemma}

\begin{proof}
Write $\mathcal L=-\Div(aD)$ and, for $s\in\{-1,1\}$, set
$V_s=e^{skZ}$.  The weak chain rule gives
\begin{equation}\label{n3:existence:exponential-identity}
 \mathcal L V_s
 =\Div(skV_sb)-skV_se
       -k^2V_s\bigl(aDZ\cdot DZ+b\cdot DZ\bigr).
\end{equation}
Choose $k\ge\max\{1,2C_0/\lambda\}$ and use
$B|DZ|\le(\lambda/2)|DZ|^2+B^2/(2\lambda)$.
Since $e^{-kM}\le V_s\le e^{kM}$, we obtain
\begin{equation}\label{n3:existence:exponential-subsolution}
 \mathcal L V_s\le\Div F_s+\nu,
 \qquad |F_s|\le ke^{kM}B,
 \qquad \nu\le C\dd x+C\mu.
\end{equation}

On a ball $B_R$ compactly contained in the patch, solve
$\mathcal L w_s=\Div F_s+\nu$ with zero Dirichlet data.
The bounded vector source has a solution of supremum norm at most
$CR$: apply the $L^p$ source estimate, $p>3$, whose scaled factor
is $R^{1-3/p}$, to $\|F_s\|_p\le CR^{3/p}$.
For the positive measure source use the Green bound
$0\le\mathcal G(x,y)\le C|x-y|^{-1}$ for uniformly elliptic
measurable divergence coefficients
\cite[Section 7, Theorem (7.1) and Remark 2]{LittmanStampacchiaWeinberger1963}.
To use that interior comparison uniformly for all $x,y\in B_R$,
extend the coefficients elliptically to $B_{2R}$.  Domain
monotonicity bounds the zero-Dirichlet Green function on $B_R$
by the one on $B_{2R}$, where $B_R$ is a fixed compact interior
region after rescaling.  This gives the stated uniform bound.
Integration over concentric annuli gives
\begin{equation}\label{n3:existence:green-morrey-correction}
 \sup_{x\in B_R}\int_{B_R}\mathcal G(x,y)\dd\nu(y)
 \le C\int_0^{2R}\frac{\nu(B_t(x))}{t^2}\dd t
       +CR^{-1}\nu(B_{2R}(x))
 \le C(R^2+R^\beta).
\end{equation}
One may first approximate the measure by smooth positive densities.
The uniform potential bound and the energy identity
$\int aDw\cdot Dw=\int w\dd\nu$ give a bounded sequence in
$W^{1,2}_0$ and hence the required weak solution.  Thus, with
$\gamma=\min\{1,\beta\}>0$,
\begin{equation}\label{n3:existence:correction-size}
 \|w_s\|_\infty\le\varepsilon_R,
 \qquad \varepsilon_R=CR^\gamma.
\end{equation}

Let $M_R=\sup_{B_R}Z$ and $m_R=\inf_{B_R}Z$.  The functions
\begin{equation}\label{n3:existence:exponential-deficits}
 H_+=e^{kM_R}-e^{kZ}+w_++\varepsilon_R,
 \qquad
 H_-=e^{-km_R}-e^{-kZ}+w_-+\varepsilon_R
\end{equation}
are nonnegative and satisfy $\mathcal L H_\pm\ge0$.
Weak Harnack controls their $L^q$ means on $B_{R/2}$ by their
infima on $B_{R/4}$.  Each exponential deficit is comparable,
with constants depending on $kM$, to $M_R-Z$ or $Z-m_R$,
respectively.  At least one of the sets
\[
 \{Z\le(M_R+m_R)/2\}\cap B_{R/2},\qquad
 \{Z\ge(M_R+m_R)/2\}\cap B_{R/2}
\]
has at least half the measure of $B_{R/2}$.  Apply weak Harnack to
the corresponding deficit.  It follows that either the upper
supremum decreases, or the lower infimum increases, by at least
$c(M_R-m_R)-C\varepsilon_R$ on $B_{R/4}$.  Consequently
\begin{equation}\label{n3:existence:oscillation-contraction}
 \operatorname{osc}_{B_{R/4}}Z
 \le(1-c)\operatorname{osc}_{B_R}Z+CR^\gamma.
\end{equation}
Iteration proves a H\"older estimate with any sufficiently small
positive exponent below both $\gamma$ and
$-\log(1-c)/\log4$.

For the boundary estimate, flatten a face and put
$S=\operatorname{diag}(1,1,-1)$.  An even extension of $Z$ uses
$a_-=Sa_+S$ and $b_-=Sb_+$.  Its flux has a plane jump whose
density is twice the prescribed normal flux, hence bounded.
An odd extension after subtracting a constant Dirichlet value uses
$a_-=Sa_+S$, $b_-=-Sb_+$, and the sign-reflected bulk source.
Its normal flux is continuous.  The transformed metrics and volume
densities are absorbed into the displayed coefficients.  A bounded
plane density has ball mass $O(r^2)$, so it is admissible with
exponent one.  This proves the asserted boundary version.
\end{proof}


\begin{proposition}[Classical bounds for bounded-variable solutions]
\label{n3:existence:classical-bounds}
Fix the deformation parameters and assume the bounds in
Proposition~\ref{n3:apriori:bounds}.  Every smooth solution in the
homotopy of Proposition~\ref{n3:deformation:homotopy}, with
$t\ge-\epsilon$, has uniform local classical estimates of every
finite order, up to the inner and outer faces.  At fixed $n$ the
constants are independent of the homotopy parameter and of all
sufficiently large outer radii $R$.  The local patches can be chosen
with a uniform positive radius throughout the end.
\end{proposition}

\begin{proof}
Start with the bounds for $f,Z,t,\sigma$ from
Proposition~\ref{n3:apriori:bounds}.
At fixed parameters they bound $u,l,D$ above and away from zero,
and give a uniform $\chi_0>0$.  After dividing the scalar trace
equation by $l/\sqrt D$, its right side is
\begin{equation}\label{n3:existence:normalized-f-source}
 G_f=\frac{\sqrt D}{l}
         \bigl(F_s-\tr_{A_\chi}K\bigr).
\end{equation}
It is bounded throughout the homotopy, including its added collar
terms.  Lemma~\ref{n3:existence:gradient} therefore gives a uniform
$C^{1,\alpha}$ estimate for $f$ and
\begin{equation}\label{n3:existence:f-morrey-application}
 \int_{B_r}|\Hess f|^2\le Cr^{1+2\alpha}
\end{equation}
on the uniform patches, including boundary patches where $f$ is
constant.

After multiplying by the Riemannian volume density, the second
equation has coordinate principal coefficients
$a^{ij}=4\sqrt{\det g}\,uA_\chi^{ij}$, a bounded flux vector $b$, and a source
bounded in absolute value by
\begin{equation}\label{n3:existence:Z-quadratic-source}
 C\bigl(1+|DZ|^2+|D^2f|^2\bigr).
\end{equation}
Indeed $t$ is a smooth function of $x,Z,Df$ on the bounded variable
range.  Thus $Dt$ is a bounded linear combination of $DZ,D^2f$
and fixed smooth terms; the expression $\mathcal T$ is quadratic
in these quantities.  All other source terms have lower order.
The same volume density is included in $b,e$; it is smooth and bounded
above and away from zero on the chosen patches, so it preserves these
bounds and uniform ellipticity.  The prescribed conormal flux at the inner boundary is
bounded.  At the outer face $Z=0$.
Apply Lemma~\ref{n3:existence:natural-growth}, with
$\mu=C|D^2f|^2\dd x$ and the bounded reflected plane density.
Equation~\eqref{n3:existence:f-morrey-application} supplies its measure
growth hypothesis.  We obtain a uniform positive H\"older exponent
for $Z$.

Now $Z$ and $Df$ are H\"older, so all coefficients and the right
side of the scalar $f$-equation are H\"older on their bounded
variable range.  Interior and constant-Dirichlet boundary Schauder
estimates give $f\in C^{2,\theta}$ for a uniform $\theta>0$.
Expand the $Z$-equation in nondivergence form.  Its principal
coefficients are uniformly elliptic and $C^{0,\theta}$, while its
right side is bounded by $C(1+|DZ|^2)$.
On an inner face $Df$ is normal.  Therefore $A_\chi$ has zero
normal--tangential components there and normal eigenvalue $\chi$.
Its boundary equation solves for
\begin{equation}\label{n3:existence:ordinary-normal-condition}
 \partial_\nu Z=\mathcal H_s(x,Z,f,Df),
\end{equation}
where $\mathcal H_s$ is smooth on a fixed bounded variable range.
Its derivatives with respect to $x$ and its displayed scalar
arguments are bounded at fixed parameters.  The outer condition
remains constant Dirichlet.  These are different, disjoint smooth
boundary components; there is no change of boundary condition along
an edge or a corner.

It remains to absorb the quadratic gradient growth. Fix
$p>3$.  The linear local $W^{2,p}$ estimate for the leading
nondivergence operator with Equation~\eqref{n3:existence:ordinary-normal-condition}
or outer Dirichlet data has the form, on fixed nested patches
$U_x\Subset V_x$,
\begin{equation}\label{n3:existence:Wp-pre-absorption}
 \|Z\|_{W^{2,p}(U_x)}
 \le C\left(1+\|DZ\|_{L^{2p}(V_x)}^2
                 +\|Z\|_{W^{1,p}(V_x)}\right).
\end{equation}
Here the boundary norm is controlled by
\begin{equation}\label{n3:existence:boundary-trace-composition}
 \|\mathcal H_s(x,Z,f,Df)\|_{W^{1-1/p,p}(\partial V_x)}
 \le C\bigl(1+\|Z\|_{W^{1,p}(V_x^+)}\bigr),
\end{equation}
using trace and composition after extending its smooth expression
through the collar.  The derivatives of $f,Df$ entering this
extension are already bounded.  The leading matrix has a fixed
H\"older modulus, so its small oscillation on sufficiently small
patches permits the usual freezing and absorption in the linear
estimate.  The normal boundary operator is uniformly oblique and
satisfies the complementing condition.  These are the classical
local elliptic boundary estimates
\cite{AgmonDouglisNirenberg1959,AgmonDouglisNirenberg1964}.

Let $[Z]_{C^{0,\theta}}\le H_0$.  On balls or half-balls of
radius $h$, the scaled Gagliardo--Nirenberg inequality, applied after
subtracting a constant, gives
\begin{equation}\label{n3:existence:small-oscillation-GN}
 \begin{split}
 \|DZ\|_{L^{2p}(B_h)}^2
 &\le C\operatorname{osc}_{B_h}Z\,
                    \|D^2Z\|_{L^p(B_h)}\\
 &\quad+Ch^{3/p-2}(\operatorname{osc}_{B_h}Z)^2.
 \end{split}
\end{equation}
One can use a bounded extension operator on a half-ball; its constants
are uniform in the boundary charts.  This is the interpolation
inequality of \cite[Lecture II, Equation (2.2)]{Nirenberg1959}
with its lower-order term on a bounded patch.
Cover $V_x$ by radius-$h$ patches of bounded overlap and take the
$\ell^p$ sum of Equation~\eqref{n3:existence:small-oscillation-GN}.
It follows that
\begin{equation}\label{n3:existence:global-local-GN}
 \|DZ\|_{L^{2p}(V_x)}^2
 \le CH_0h^\theta\|D^2Z\|_{L^p(V_x^+)}
           +CH_0^2h^{2\theta-2}.
\end{equation}
Also, for any $\eta>0$,
$\|Z\|_{W^{1,p}(V_x)}
 \le\eta\|D^2Z\|_{L^p(V_x^+)}+C_\eta$.
Each enlarged patch $V_x^+$ is covered by a bounded number of the
patches $U_y$, independently of $R$.  Set
$S=\sup_y\|Z\|_{W^{2,p}(U_y)}$, finite for each smooth finite-domain
solution.  Equations~\eqref{n3:existence:Wp-pre-absorption} and
\eqref{n3:existence:global-local-GN} give
\begin{equation}\label{n3:existence:uniform-local-absorption}
 S\le C(H_0h^\theta+\eta)S+C(h,\eta).
\end{equation}
Choose $h$ and then $\eta$ so the first coefficient is less than
$1/2$.  We obtain a uniform $W^{2,p}$ bound and hence a
$C^{1,1-3/p}$ bound for $Z$.

The coefficients in both equations and in the normal boundary
condition now have the regularity required for Schauder bootstrap.
First the nondivergence $Z$-equation has H\"older source, giving
$Z\in C^{2,\theta'}$; the scalar $f$-equation then gains another
derivative.  Repetition gives every finite classical derivative
supported by the smooth data.  The same procedure holds on the
uniform unit patches of the asymptotic end and on the large outer
spheres, whose normal-coordinate constants are uniform.  This
proves the asserted independence of $R$.  The reduced background
has end symbol bounds of every order by
Proposition~\ref{n3:reduction:rest-reduction}; the bootstrap here
does not assume additional derivatives of the original asymptotic
data.
\end{proof}

\subsection{The scalar Dirichlet solution operator}

The compact map requires a scalar solve for every bounded input $Z$.
Its definition must therefore work without a floor bound on the input.

\begin{lemma}[Scalar solve for unrestricted bounded inputs]
\label{n3:existence:scalar-solve}
Fix $n\ge4$, $\tau>0$, a finite smooth truncation $\Omega_R$,
and a homotopy parameter $s\in[0,4]$.  For every
$Z\in C^{1,b}(\overline{\Omega_R})$, $0<b<1$, the scalar trace
equation of Proposition~\ref{n3:deformation:homotopy}, with its
prescribed componentwise constant Dirichlet values for $f$, has a
unique solution.  It lies in $C^{3,b}$ and is smooth to the extent
allowed by the input.  The solution operator
\begin{equation}\label{n3:existence:scalar-operator}
 f=\mathcal S_s(Z)
\end{equation}
is continuous in $(s,Z)$, including across the concatenation points
of the homotopy.  On bounded subsets of $C^{1,b}$ its $C^{3,b}$
norm is uniformly bounded at these fixed parameters.  No lower
bound for $t$ is required.
\end{lemma}

\begin{proof}
The implicit definition of $t$ has derivative
$\partial Z/\partial t=1+pv/4>0$.  Thus $t$ and all the
coefficients are smooth functions of $x,Z,Df$, also at $Df=0$.
There is no dependence of the principal coefficient on the height
$f$ itself.  For bounded height the added collar terms in $F_s$
are bounded, and its derivative with respect to $h=\tau f$ is at
least one.  At $Df=0$ the equation and this height sign give
upper and lower constant barriers.  Hence every scalar solution
satisfies
\begin{equation}\label{n3:existence:scalar-height}
 |h|\le C_0,
\end{equation}
where the constant can be chosen independently of $Z$ and the
homotopy parameter.  We include the fixed boundary values in its
choice.

We first derive the coefficient bounds at an unbounded slope.
If $|Z|\le M$ and $\sigma\to\infty$, the implicit equation forces
$t\to-\infty$ uniformly in $Z$.  There $p=n$ and $l=e^{nt}$,
so its exact form is
\begin{equation}\label{n3:existence:implicit-large-slope}
 e^{8Z}=e^{8t}+e^{(2n+8)t}\sigma^2.
\end{equation}
Consequently
\begin{equation}\label{n3:existence:large-slope-asymptotics}
 \begin{split}
 t&=-\frac{\log\sigma}{n+4}+O(1),\qquad
 l\asymp\sigma^{-n/(n+4)},\\
 d&\asymp\sigma^{-8/(n+4)},\qquad
 \chi\asymp\sigma^{-8/(n+4)},\qquad
 \frac{\sqrt D}{l}\asymp\sigma.
 \end{split}
\end{equation}
The constants are uniform for $|Z|\le M$ at fixed $n$.
Since $n\ge4$, the exponent $8/(n+4)$ is at most one.  Combining
the large-slope estimates with compactness of the remaining variable
range yields
\begin{equation}\label{n3:existence:scalar-growth}
 \chi\ge\frac{c}{1+\sigma},\qquad
 |G_f|\le C(1+\sigma)
\end{equation}
whenever Equation~\eqref{n3:existence:scalar-height} holds.

We prove a global gradient estimate using this radial ellipticity.
Let $\delta$ be smaller than a fixed collar radius, half the distance
between distinct boundary components, and a normal injectivity
radius of a smooth extension of the finite domain.  On a boundary
distance collar use
\begin{equation}\label{n3:existence:logarithmic-modulus}
 \psi(r)=k^{-1}\log(1+Br),\qquad
 \psi''=-k(\psi')^2.
\end{equation}
The upper and lower barriers are the assigned face value plus or
minus $\psi$.  On these tests the axis is normal to the distance
leaves, so their principal trace is
$\chi\psi''+\psi'\Delta r$.  At slopes at least one,
Equation~\eqref{n3:existence:scalar-growth} gives
\begin{equation}\label{n3:existence:logarithmic-dominance}
 -k\chi(\psi')^2\le-\tfrac12ck\psi'.
\end{equation}
Choose $k$ large enough to dominate all the bounded geometry terms
and the linear growth in Equation~\eqref{n3:existence:scalar-growth}.
Next shrink $\delta$ so that $1/(2k\delta)$ exceeds any fixed
large-slope threshold used here.  Finally choose $B$ large enough
that $B\delta\ge1$ and $\psi(\delta)$ exceeds twice the height
oscillation.  Then $\psi'\ge1/(2k\delta)$ on this collar and
the two barriers compare by the height monotonicity.  This gives
a uniform boundary gradient bound.

For the interior estimate, increase $B$ if necessary and compare
two points using
\begin{equation}\label{n3:existence:doubling-function}
 f(x)-f(y)-\psi(\operatorname{dist}(x,y)),
 \qquad 0<\operatorname{dist}(x,y)<\delta.
\end{equation}
The distance is taken in the fixed smooth extension.  The
boundary barriers exclude a positive maximum with one endpoint on
the boundary.  The inequality $\psi(\delta)>2\operatorname{osc}f$
excludes the other boundary of the two-point domain.  At a positive
interior maximum let $r=\operatorname{dist}(x,y)$ and
$q=\psi'(r)$.  The two gradients have length $q$ and are parallel
transports along the connecting geodesic.  Vary both endpoints
simultaneously in parallel transverse directions.  The second
variation formula for the short geodesic yields
\begin{equation}\label{n3:existence:two-point-transverse}
 \tr_{e_x^\perp}\Hess f(x)
 -\tr_{e_y^\perp}\Hess f(y)\le Cqr.
\end{equation}
Varying each endpoint separately along that geodesic gives
$f_{;e_xe_x}(x)\le\psi''(r)$ and
$f_{;e_ye_y}(y)\ge-\psi''(r)$.  Subtracting the two scalar
equations therefore gives
\begin{equation}\label{n3:existence:two-point-contradiction}
 -2C(1+q)
 \le G_f(x)-G_f(y)
 \le Cqr+(\chi_x+\chi_y)\psi''(r).
\end{equation}
By Equations~\eqref{n3:existence:scalar-growth} and
\eqref{n3:existence:logarithmic-modulus}, the last term is at most
$-ckq$ at the large comparison slopes.  Our choice of $k$ gives
a contradiction.  Thus Equation~\eqref{n3:existence:doubling-function}
is nonpositive.  Reversing $x,y$ proves the same bound for the
absolute difference and, on letting $y\to x$, a global Lipschitz
bound for $f$.  This argument used no continuity modulus for
$\chi_x-\chi_y$; only the common transverse eigenvalues and the
radial lower bound were needed.

For existence interpolate the normalized equation with
$\Delta f=\tau f$, retaining the given boundary values.  Explicitly,
for $a\in[0,1]$ take principal matrix
$(1-a)I+aA_\chi$ and right side
$(1-a)\tau f+aG_f$.  The height sign, the linear gradient growth,
and Equation~\eqref{n3:existence:scalar-growth} persist uniformly in
$a$.  The height and logarithmic gradient bounds just proved are
therefore uniform.  After these bounds the interpolating equations
are uniformly elliptic.  Their matrices still have the axial form,
with axial eigenvalue $(1-a)+a\chi$, so
Lemma~\ref{n3:existence:gradient} applies.  Schauder estimates then
give a uniform $C^{2,\theta}$ bound at fixed input $Z$.
The coefficients and right side are now Lipschitz in position on
bounded input sets, which permits their $C^{0,b}$ estimates and
an upgrade to $C^{2,b}$.  Differentiating once, or using the next
Schauder estimate, yields $f\in C^{3,b}$ because $Z\in C^{1,b}$.

The scalar linearization has the form
$a^{ij}D_{ij}\varphi+b^iD_i\varphi-c\varphi$ with
$c\ge c_1\tau>0$ on the bounded variable range.  Indeed the
principal matrix is independent of the height, while
$\partial_fG_f\ge\tau\sqrt D/l>0$.
The Dirichlet maximum principle gives uniqueness and a zero kernel;
the linear elliptic Dirichlet theory gives an inverse.  The
implicit-function theorem gives openness of the interpolation
parameter set.  The uniform classical bounds and compactness give
closedness, starting from the invertible Laplace endpoint.  This
proves existence.  The same maximum principle applied to the
difference of two solutions proves uniqueness of the nonlinear
solution.  The implicit-function theorem, or compactness and
uniqueness, gives continuous dependence on $Z$ and on the
homotopy parameter.  At the finitely many concatenation points the
scalar equations agree, so the same continuity holds there.
\end{proof}


\subsection{A compact equation and the order of parameter choices}

Recall the parameter $s\in[0,4]$ of
Proposition~\ref{n3:deformation:homotopy}: $s=0$ is the desired system,
$s=1$ has zero drift, $s=2$ has the full collar modification, and
$s=3$ has $f=0$. On $[3,4]$ the scalar source and inner flux are scaled
to zero; at $s=4$ the scalar equation still has the unique solution
$f=0$. We will define a compact map along this path and show that its
fixed points remain in the admissible set.

\begin{theorem}[Existence on finite truncations]
\label{n3:existence:deformation}
For the prepared domain of Proposition~\ref{n3:domains:prepared-domain}
and every sufficiently small fixed $\epsilon>0$, there is an
integer $n_0\ge4$ and, for each $n\ge n_0$, a radius
$R_{\min}(n)$, with $R_{\min}(n)\to\infty$, such that the original
deformation system has a smooth solution on every $\Omega_R$ with
$R\ge R_{\min}(n)$.  It satisfies $t>-\epsilon$ on the closure.
At fixed $n$, these solutions have the local bounds of
Proposition~\ref{n3:existence:classical-bounds}, uniformly as
$R\to\infty$.  In particular, any sequence $R_j\to\infty$ of such radii has a
subsequence converging smoothly on compact subsets of
$\overline\Omega$ to a solution of the original equations with
$t\ge-\epsilon$.
\end{theorem}

\begin{proof}
We first make the radius convention precise.  Fix the prepared
data, all thresholds and collars, and $\epsilon$.  Increase $n_0$
so that
\begin{equation}\label{n3:existence:n-choice}
 n_0\ge\max\{4,\lceil2/\epsilon\rceil\}
\end{equation}
and so that the floor test and all estimates of
Proposition~\ref{n3:apriori:bounds} apply for $n\ge n_0$.
Let $R_{\mathrm{ap}}(n)$ be a sufficient lower radius for those
estimates, and let $R_{\mathrm{geo}}$ contain all compact boundary
faces and supports used in the construction.  The end barrier of
that proposition gives a constant $C_{\mathrm{out}}$, independent
of $n,R,s$, such that every admissible or floor-contact homotopy
solution has
\begin{equation}\label{n3:existence:outer-gradient}
 |\nabla f|\le
       C_{\mathrm{out}}\tau^{-1}R^{-1-\gamma}
       \quad\text{on the outer sphere}.
\end{equation}
The scalar last stage has $f=0$ and satisfies the same bound.
Choose, enlarging to a radius of a smooth truncation if necessary,
\begin{equation}\label{n3:existence:R-choice}
 \begin{split}
 R_{\min}(n)=\max\biggl\{&n,R_{\mathrm{geo}},R_{\mathrm{ap}}(n),\\
 &\left(\frac{2C_{\mathrm{out}}\ell}
               {\tau\sqrt{e^{8\epsilon}-1}}
        \right)^{1/(1+\gamma)}\biggr\}.
 \end{split}
\end{equation}
At $Z=0$, the value of its implicit expression at $t=-\epsilon$
is $-\epsilon+\frac18\log(1+\ell^2\sigma^2)$.
Equations~\eqref{n3:existence:outer-gradient} and
\eqref{n3:existence:R-choice} make this strictly negative.  Strict
monotonicity of that expression in $t$ therefore excludes an outer
floor contact.  The term $n$ in Equation~\eqref{n3:existence:R-choice}
ensures $R_{\min}(n)\to\infty$.

The sublevel and collar estimates apply to arbitrary smooth
admissible or floor-contact solution sequences with $n\to\infty$
and expanding outer cutoffs, uniformly over the compact ranges of
the homotopy parameters.  If the needed conclusion failed for
arbitrarily large $n$ with $R\ge R_{\min}(n)$, one could select
such a violating sequence.  The chosen radius convention would
make it precisely an expanding sequence excluded by those proofs.
Thus a single $n_0$, followed by the radius choices above, suffices
for every homotopy parameter.  No upper bound for $Z$ is inserted
in the earlier contact estimates: there $\sigma\asymp\tau^{-1}$
and $l\ge\ell$ alone give $d\lesssim(\tau/\ell)^2$.

Fix now $n\ge n_0$ and $R\ge R_{\min}(n)$ for the rest of the
degree argument.  Let
\begin{equation}\label{n3:existence:banach-space}
 X_b=\{Z\in C^{1,b}(\overline{\Omega_R}):
                       Z=0\text{ on the outer face}\},
 \qquad 0<b<1.
\end{equation}
For an input $Z\in X_b$, first solve
$f=\mathcal S_s(Z)$ by Lemma~\ref{n3:existence:scalar-solve}.
Compute $t,u,A_\chi,\mathcal T$ and the full homotopy source
from this input pair.  Write $B_s$ for the frozen drift summand in
the flux, $E_s$ for the full right side of its divergence equation,
and $q_s$ for the prescribed inner flux.  Thus before the final
scaling $B_s=\lambda_suA_\chi K(w,\cdot)$ and $E_s=u\Xi$;
on the final interval $E_s$ and $q_s$ have the simultaneous scalar
factor prescribed in Proposition~\ref{n3:deformation:homotopy}.
Define $\mathcal K_s(Z)=\widetilde Z$ by the linear mixed problem
\begin{equation}\label{n3:existence:compact-map-equation}
 \begin{cases}
  \Div(4uA_\chi\nabla\widetilde Z+B_s)=E_s
                           &\text{in }\Omega_R,\\
  (4uA_\chi\nabla\widetilde Z+B_s)_\nu=q_s
                           &\text{on }B,\\
  \widetilde Z=0           &\text{on the outer face}.
 \end{cases}
\end{equation}
The normal $\nu$ in the second line points into $\Omega_R$; the
variational formulation uses its negative as the integration
normal.  This changes the sign of the boundary functional, not
coercivity.  On a bounded input set, the scalar gradient estimate
bounds the coefficients above and away from zero.  The nonempty
outer Dirichlet face gives a Poincar\'e inequality, so the linear
problem is uniquely solvable by its coercive bilinear form.

On bounded input sets in $C^{1,b}$,
Lemma~\ref{n3:existence:scalar-solve} bounds $f$ in $C^{3,b}$.
Hence the leading coefficient and drift in
Equation~\eqref{n3:existence:compact-map-equation} are $C^{1,b}$,
the bulk source is $C^{0,b}$, and the prescribed boundary data are
$C^{1,b}$.  The derivatives of $t(x,Z,Df)$ involve only $DZ$
and $D^2f$, so this assertion includes the quadratic expression
$\mathcal T$.  Linear regularity gives a bounded output set in
$C^{2,b}$.  The inclusion $C^{2,b}\hookrightarrow C^{1,b}$ is
compact on the finite domain.  Continuous dependence of the scalar
solve and the linear problem makes $(s,Z)\mapsto\mathcal K_s(Z)$
a continuous compact homotopy on bounded sets.  These assertions
hold also at the concatenation parameters because the defining
equations coincide there.  A fixed point solves the coupled system;
repeated Schauder estimates make it smooth.

Admissibility is evaluated after the scalar solve.  Define the open
subset of the parameter--input product
\begin{equation}\label{n3:existence:admissible-open-set}
 \mathcal O=\{(s,Z)\in[0,4]\times X_b:
       \min_{\overline{\Omega_R}}t_s[Z]>-\epsilon,
       \ \|Z\|_{C^{1,b}}<M_1\}.
\end{equation}
It is relatively open because the scalar solve and the implicit
function defining $t$ depend continuously on $(s,Z)$.
Propositions~\ref{n3:apriori:bounds} and
\ref{n3:existence:classical-bounds} give a uniform bound in $C^{1,b}$
for all fixed points satisfying $t\ge-\epsilon$.  If initially
needed, use their higher classical bounds to reach the chosen
exponent $b$.  Choose $M_1$ larger than this bound.  No fixed point
lies on the norm boundary of $\mathcal O$.  No fixed point lies on
its floor boundary either: outer contacts were excluded above and
all remaining contacts are excluded by
Proposition~\ref{n3:deformation:floor}.

To apply degree on the moving admissible set, we need compactness of
its fixed points.
The set of fixed points in $\overline{\mathcal O}$ is compact:
their $X_b$ norms are bounded, the parameter interval is compact,
and $Z=\mathcal K_s(Z)$ with a compact homotopy.  Any limit has
$t\ge-\epsilon$ by continuity; it is a smooth fixed point and
cannot lie on the floor.  Thus these fixed points are compactly
separated from the excluded boundary.  The excision form of
Leray--Schauder homotopy invariance applies to the sections
$\mathcal O_s$; equivalently, cover the compact fixed-point set
by finitely many product neighborhoods contained in
$\mathcal O$ and use ordinary homotopy invariance and excision on
the resulting finite parameter subdivision.  This is the compact
perturbation-of-identity degree \cite{LeraySchauder1934}.

At $s=4$, the scalar equation has the unique solution $f=0$.
The drift vanishes, and both $E_4$ and $q_4$ vanish for every input.
Equation~\eqref{n3:existence:compact-map-equation} then has output
identically zero.  Moreover, $t_4[0]=0>-\epsilon$, so
$0\in\mathcal O_4$ and
\begin{equation}\label{n3:existence:degree-one}
 \deg(I-\mathcal K_4,\mathcal O_4,0)=1.
\end{equation}
Homotopy invariance gives the same degree at $s=0$.  There is
therefore a fixed point in $\mathcal O_0$, which is the desired
smooth solution with $t>-\epsilon$.

The bounds of Proposition~\ref{n3:existence:classical-bounds} are
independent of $R$ at fixed $n$.  A diagonal subsequence on a
compact exhaustion, using boundary charts at the fixed inner
faces, converges smoothly to a solution on $\overline\Omega$.
The floor inequality passes to $t\ge-\epsilon$. The next section
establishes the end normalization, decay, and flux for these limits.
\end{proof}
